% !TEX program = xelatex
% main.tex — Proyecto final: matriz SABSA para pagos inmediatos Bre-B.
% (la primera línea le dice a Texel y a otros editores que use XeLaTeX: el tema lo necesita)
% Compilar con XeLaTeX:  latexmk main.tex   (ver latexmkrc)
\documentclass[aspectratio=169,10pt]{beamer}
\usetheme[nosectionpages,accent=DA0921]{terminal}  % rojo oficial UR
\usepackage[spanish,es-noshorthands]{babel}

% caso.tex — datos del caso que cambian si se cambia de empresa.
% El contenido completo de la matriz vive en ../matriz/*.md (fuente del Excel);
% las slides lo resumen a mano.
\newcommand{\empresa}{NovaPay S.A.}
\newcommand{\empresaslug}{novapay}  % en los prompts de shell
\newcommand{\empresatipo}{compañía de financiamiento digital vigilada por la SFC}
\newcommand{\area}{Gerencia de Pagos Inmediatos}
\newcommand{\proyecto}{Vinculación a Bre-B como participante}
\newcommand{\autores}{Sergio Andrés Caceres Díaz\texorpdfstring{\\}{, } Over Haider Castrillón Valencia\texorpdfstring{\\}{, } Michael Steven Pinilla Mateus}
\newcommand{\autorescorto}{Caceres, Castrillón, Pinilla}  % pie de página


\title{Arquitectura de seguridad SABSA para pagos inmediatos Bre-B}
\subtitle{\empresa{} // Proyecto: \proyecto}
\author[\autorescorto]{\autores}
\institute{Seguridad para Arquitectura Empresarial -- Universidad del Rosario}
\date{Noviembre 2026}

\terminalwindow{sabsa-breb.pdf}
\terminalteamlabel{equipo}
\terminallogo{\includegraphics[height=2.6ex]{pics/ur-logo-white}}
\terminaltitlelogo{\includegraphics[height=1.25cm]{pics/ur-logo}}
\terminalfootprefix{sabsa}
\terminalsectionprompt{cd capa\_}
\terminalshell{\texttt{\$ ./sabsa -{}-empresa \empresaslug{} -{}-proyecto bre-b\\
  > capas contextual..component ...... ok\\
  > matriz SABSA (xlsx) .............. ok}}

\begin{document}

\begin{frame}[plain,noframenumbering]
  \titlepage
\end{frame}

% ------------------------------------------------------------------ 01
\section{Misión}
\begin{frame}{Misión y visión de la empresa}
  \prompt{cat /etc/\empresaslug/\{mision,vision\}}
  \vskip6pt
  \begin{columns}[T]
    \begin{column}{0.56\textwidth}
      \begin{tbox}{> mision.txt}
        Facilitar que personas y comercios en Colombia muevan su dinero de
        forma \acc{inmediata}, \acc{segura} y \acc{sin costo} desde el celular,
        protegiendo la \acc{confianza} de los clientes.
      \end{tbox}
    \end{column}
    \begin{column}{0.4\textwidth}
      \begin{tbox}{> vision.txt}
        Ser en 2030 la billetera digital preferida por quienes hoy usan
        efectivo, reconocida por la seguridad de sus pagos.
      \end{tbox}
    \end{column}
  \end{columns}
  \vskip14pt
  \begin{tbox}[colframe=terminalaccent, colbacktitle=terminalaccent]{> objetivo\_estrategico\_2028}
    \large Crecer 25\,\% en clientes activos a diciembre de 2028 \acc{sin que el
    fraude pase de 5 pb} del valor transado.
  \end{tbox}
  \vskip10pt
  {\footnotesize\textbf{\empresa}: \empresatipo.\par}
\end{frame}

% ------------------------------------------------------------------ 02
\section{Proyecto}
\begin{frame}{Objetivo del área y del proyecto}
  \prompt{sabsa init -{}-proyecto bre-b -{}-dueño "\area"}
  \vskip4pt
  \begin{tbox}{\texttt{> \area: objetivo.txt}}
    Vincular a \empresa{} a \acc{Bre-B} como participante (vía un sistema de pago certificado) y lanzar en la
    app las transferencias por llaves y los pagos QR (\acc{24/7}, \acc{< 20 s}),
    con la seguridad incluida desde el diseño.
  \end{tbox}
  \vskip8pt
  % diagrams/alineacion.tex — por qué existe el proyecto: de la misión a la arquitectura, y sus hitos
% con la actividad de seguridad de cada fase (para que la seguridad no se quede por fuera del cronograma).
\begin{tikzpicture}[x=1cm, y=1cm]
  % cadena de alineación (la arquitectura alinea el objetivo de negocio con TI)
  \begin{scope}[every node/.style={terminal box, text width=2.3cm, minimum height=1.05cm}]
    \node (mis) at (1.45,0) {\textbf{MISIÓN}\\pagos inmediatos,\\seguros y gratis};
    \node (obj) at (5.15,0) {\textbf{OBJETIVO 2028}\\+25\,\% clientes\\fraude < 5 pb};
    \node (pro) at (8.85,0) {\textbf{PROYECTO}\\entrar a Bre-B\\llaves y QR};
    \node[terminal accent, line width=1pt] (arq) at (12.55,0)
      {\textbf{ARQUITECTURA}\\SABSA: 5 capas\\negocio → producto};
  \end{scope}
  \draw[terminal arrow] (mis) -- node[above, font=\tiny] {orienta} (obj);
  \draw[terminal arrow] (obj) -- node[above, font=\tiny] {justifica} (pro);
  \draw[terminal arrow, draw=terminalaccent] (pro) -- node[above, font=\tiny, text=terminalaccent] {se protege} (arq);

  % hitos del proyecto (fase arriba en negro, actividad de seguridad abajo en rojo)
  \draw[line width=0.8pt] (0.2,-1.55) -- (13.8,-1.55);
  \foreach \x/\q/\txt/\seg [count=\k] in {1.45/Q1 2027/vinculación con BanRep/{licencias firmadas y\\análisis de riesgos},
                               5.15/Q2 2027/desarrollo y certificación/{pruebas de seguridad\\y pentest},
                               8.85/Q3 2027/marcha blanca/{SOC 24/7\\y playbooks}} {
    \fill (\x,-1.55) circle (2.2pt);
    \node[terminal tag, font=\tiny\bfseries, above=4pt] at (\x,-1.55) {\q};
    \node[font=\tiny, below=4pt, align=center] (t\k) at (\x,-1.55) {\txt};
    \node[font=\tiny, text=terminalaccent, align=center, below=1pt of t\k] {\seg};
  }
  \fill[terminalaccent] (12.55,-1.55) circle (2.6pt);
  \node[terminal tag, font=\tiny\bfseries, fill=terminalaccent, above=4pt] at (12.55,-1.55) {Q4 2027};
  \node[font=\tiny\bfseries, text=terminalaccent, below=4pt, align=center] (tq4) at (12.55,-1.55) {lanzamiento al público};
  \node[font=\tiny, text=terminalaccent, align=center, below=1pt of tq4] {prueba de recuperación y\\aprobación de seguridad};
  \node[font=\tiny, anchor=north west, text=black!70] at (0.2,-2.75) {en rojo: la actividad de seguridad de cada fase};
\end{tikzpicture}

\end{frame}

% ------------------------------------------------------------------ 03
\section{Contextual}
\begin{frame}{Contextual \textmd{// requisitos de negocio}}
  \prompt{sabsa --layer contextual --view business --owner "Junta + Presidencia"}
  \vskip4pt
  \renewcommand{\sabsacellfont}{\scriptsize}
  \begin{sabsagrid}
    \sabsacell{WHAT}{activos y objetivos}{\begin{cellitems}
      \item \acc{Confianza}: El activo principal.
      \item Pagos a cualquier entidad, al instante y gratis.
      \item Fondos y datos \acc{confidenciales, íntegros y disponibles}.
    \end{cellitems}}
    \sabsacell{WHY}{oport. / amenazas}{\begin{cellitems}
      \item \acc{+} 110\,M de llaves y 66\,\% usa sobre todo efectivo.
      \item \acc{+} Comercios QR y marca segura.
      \item $-$ Fraude irreversible, caídas, sanciones y fugas.
    \end{cellitems}}
    \sabsacell{HOW}{procesos}{\begin{cellitems}
      \item Vinculación, desarrollo, pruebas y aceptación.
      \item Pagos 24/7, fraude, cuadre y servicio.
      \item Registro de cada transacción para auditarla.
    \end{cellitems}}
    \sabsacell{WHO}{sponsors y extendida}{\begin{cellitems}
      \item Accionistas, Junta, Presidencia y VP.
      \item Dueño: \area.
      \item \acc{CISO transversal} (reporta a Presidencia).
      \item Clientes, proveedores, BanRep, SFC y SIC.
    \end{cellitems}}
    \sabsacell{WHERE}{sitios}{\begin{cellitems}
      \item Sede en Bogotá, servicio nacional.
      \item Nube en EE.\,UU. (principal y alterno).
      \item Conexión a Bre-B en Bogotá.
      \item Rige Colombia (SFC, BanRep y SIC).
    \end{cellitems}}
    \sabsacell{WHEN}{tiempos objetivo}{\begin{cellitems}
      \item \acc{Lanzamiento: Q4 2027}.
      \item +25\,\% clientes a dic-2028, fraude < 5\,pb.
      \item 24/7, 99.9\,\% al mes y < 20\,s.
      \item Picos: quincenas, primas y diciembre.
    \end{cellitems}}
  \end{sabsagrid}
  \src{BanRep CRE DSP-465 (v. 31-mar-2026), SFC CBJ (CE 006/2025), BanRep, Bre-B un año (6-oct-2026)}
\end{frame}

% ------------------------------------------------------------------ 04
\section{Conceptual}
\begin{frame}{Conceptual \textmd{// estrategias de seguridad}}
  \prompt{sabsa --layer conceptual --view architect --owner "CISO + directores de área"}
  \vskip4pt
  % diagrams/conceptual.tex — capa Conceptual: atributos (en orden de importancia), dominios de
% seguridad y estrategias, y debajo objetivos de control, dueños, marco de riesgo y normas.
\begin{tikzpicture}[x=1cm, y=1cm]
  % perfil de atributos SABSA, primero los más importantes (como en la matriz)
  \tikzset{attr/.style={terminal tag, font=\tiny\bfseries, inner xsep=3pt}}
  \node[font=\scriptsize\bfseries, anchor=west] (lbl) at (0,0.32) {atributos};
  \node[attr, right=5pt of lbl, fill=terminalaccent] (a1) {RESISTENTE AL FRAUDE};
  \node[attr, right=2.5pt of a1] (a2) {DISPONIBLE};
  \node[attr, right=2.5pt of a2] (a3) {OPORTUNO};
  \node[attr, right=2.5pt of a3] (a4) {AUTENTICADO};
  \node[attr, right=2.5pt of a4] (a5) {ÍNTEGRO};
  \node[attr, right=2.5pt of a5] (a6) {CONFIDENCIAL};
  \node[attr, right=2.5pt of a6] (a7) {NO REPUDIABLE};
  \node[attr, right=2.5pt of a7] (a8) {CUMPLIDO};

  % dominios de seguridad
  \begin{scope}[every node/.style={terminal box, text width=2.2cm, minimum height=1.05cm}]
    \node (cli) at (1.45,-1.35) {\textbf{CLIENTE}\\canal y equipo\\del cliente};
    \node (dmz) at (5.0,-1.35) {\textbf{EXPOSICIÓN}\\frontera pública\\(única entrada)};
    \node (core) at (8.55,-1.35) {\textbf{CORE DE PAGOS}\\llaves, envío,\\recepción y QR};
    \node (dat) at (12.1,-1.35) {\textbf{INFO. SENSIBLE}\\clientes, pagos\\y secretos};
    \node[minimum height=0.85cm] (ges) at (3.3,-3.1) {\textbf{GESTIÓN}\\administración\\y monitoreo};
    \node[terminal accent, line width=1pt, minimum height=0.85cm] (ibk) at (10.3,-3.1) {\textbf{INTERBANCARIO}\\red Bre-B (BanRep)};
  \end{scope}

  \draw[terminal arrow] (cli) -- node[above, font=\tiny] {cifrado} (dmz);
  \draw[terminal arrow] (dmz) -- node[above, font=\tiny] {autoriz.} (core);
  \draw[terminal arrow] (core) -- node[above, font=\tiny] {mínimo} (dat);
  \draw[terminal arrow, draw=terminalaccent] (core.south) |- node[pos=0.3, left, align=right, font=\tiny\bfseries, text=terminalaccent] {canal dedicado\\y autenticado} (ibk.west);
  % Gestión entra por un bus propio a media altura, lejos de las sombras de arriba
  \coordinate (bus) at ($(ges.north) + (0,0.4)$);
  \draw[dashed, line width=0.8pt] (ges.north) -- (bus);
  \draw[terminal arrow, dashed] (bus) -| (dmz.south);
  \draw[terminal arrow, dashed] (bus) -| (cli.south);
  \node[font=\tiny, anchor=west] at ([xshift=3pt]$(ges.north)!0.5!(bus)$) {solo por acceso controlado};

  % marco Zero Trust
  \node[draw, dashed, line width=0.6pt, inner xsep=7pt, inner ysep=12pt, fit=(cli)(dat)(ges)(ibk)] (zt) {};
  \node[terminal tag, anchor=west, fill=terminalaccent] at ([xshift=8pt]zt.north west)
    {ZERO TRUST + DEFENSA EN PROFUNDIDAD};
  \node[terminal chip, font=\tiny, anchor=east] at ([xshift=-8pt]zt.north east)
    {+ resiliencia, fraude como proceso continuo, DevSecOps, privacidad e IA gobernada};

  % objetivos de control, dueños, marco de riesgo y normas
  \tikzset{row/.style={terminal chip, font=\tiny, anchor=north west, align=left}}
  \node[row] (ctl) at ([yshift=-6pt]zt.south west)
    {\textbf{control}\ \ autenticación fuerte, scoring de fraude, cifrado en todo el ciclo, recuperación, mínimo privilegio, seguridad desde el diseño};
  \node[row] (own) at ([yshift=-3pt]ctl.south west)
    {\textbf{dueños}\ \ VP de Pagos (proceso), CISO (riesgo ciber), Gte. de Fraude (fraude)};
  \node[row, terminal accent, right=4pt of own.north east, anchor=north west] (rsk)
    {\textbf{riesgo}\ \ ISO 31000 + ISO/IEC 27005 con el SIAR, en cada fase};
  \node[row] at ([yshift=-3pt]own.south west)
    {\textbf{normas}\ \ ISO 27001, NIST CSF 2.0, ISO 42001 (IA), CSA CCM, SFC (CBJ y SIAR), BanRep (DSP-465, 470 y 471), Ley 1581};
\end{tikzpicture}

\end{frame}

% ------------------------------------------------------------------ 05
\section{Logical}
\begin{frame}{Logical \textmd{// servicios de seguridad}}
  \prompt{sabsa --layer logical --view designer --owner "CISO + Arquitecto de Seguridad"}
  \vskip4pt
  % diagrams/logical.tex — capa Logical: servicios de seguridad sobre el recorrido de un pago Bre-B.
\begin{tikzpicture}[x=1cm, y=1cm]
  % pasos del pago: (nombre, x, quién actúa, servicio de seguridad)
  \foreach \n/\x/\paso/\quien/\svc in {
      1/1.1/AUTENTICAR/cliente/{autenticación\\fuerte},
      2/3.46/RESOLVER LLAVE/{NovaPay → DICE}/{resolución con\\privacidad},
      3/5.82/SCORING DE FRAUDE/NovaPay/{servicio\\antifraude},
      4/8.18/CONFIRMAR Y FIRMAR/cliente/{firma y\\no repudio},
      5/10.54/ENVIAR Y LIQUIDAR/{NovaPay → BanRep}/{mensajería inter-\\bancaria segura},
      6/12.9/NOTIFICAR/ambas partes/{notificación\\y auditoría}} {
    \node[terminal box, text width=1.75cm, minimum height=1.05cm, font=\scriptsize\bfseries]
      (p\n) at (\x,0) {[\n]\\\paso};
    \node[font=\tiny, above=3pt of p\n] {\quien};
    \node[terminal chip, font=\tiny, text width=1.75cm, align=center, below=7pt of p\n] (s\n) {\svc};
  }
  \foreach \a/\b in {1/2,2/3,3/4,4/5,5/6} \draw[terminal arrow] (p\a) -- (p\b);

  % lo crítico en acento: el scoring y la firma (no repudio)
  \node[font=\tiny\bfseries, text=terminalaccent, below=1pt of s3] {< 300 ms};
  \node[font=\tiny\bfseries, text=terminalaccent, below=1pt of s4] {no repudio};

  % tiempo total
  \draw[decorate, decoration={brace, mirror, amplitude=5pt}, line width=0.8pt]
    ([yshift=-14pt]s1.south west) -- ([yshift=-14pt]s6.south east)
    node[midway, below=6pt, font=\scriptsize\bfseries, text=terminalaccent] {< 20 s extremo a extremo};

  % perfiles de privilegio y políticas (todas las etiquetas de la franja a la misma altura)
  \tikzset{strip/.style={font=\tiny, text height=1.35ex, text depth=0.4ex}}
  \tikzset{striptag/.style={terminal tag, strip, font=\tiny\bfseries, minimum width=1.45cm}}
  % activos de información (Logical WHAT, corrección de clase)
  \node[striptag, anchor=west] (act) at (0.1,-3.0) {ACTIVOS};
  \node[terminal chip, strip, right=4pt of act] (v1) {transaccional: pagos y scoring};
  \node[terminal chip, strip, right=3pt of v1] (v2) {clientes: identidad, llaves, saldos};
  \node[terminal chip, strip, right=3pt of v2, terminal accent] (v3) {llaves cripto, secretos y logs};
  \node[terminal chip, strip, right=3pt of v3] (v4) {código y APIs};

  \node[striptag, anchor=west] (perf) at (0.1,-3.6) {PERFILES};
  \node[terminal chip, strip, right=4pt of perf] (q1) {cliente: sus llaves};
  \node[terminal chip, strip, right=3pt of q1] (q2) {agente: datos enmascarados};
  \node[terminal chip, strip, right=3pt of q2] (q3) {analista: retener/liberar};
  \node[terminal chip, strip, right=3pt of q3] (q4) {ingeniero: JIT};
  \node[terminal chip, strip, right=3pt of q4] (q5) {auditor: solo lectura};

  \node[striptag, anchor=west] (pol) at (0.1,-4.2) {POLÍTICAS};
  \node[terminal chip, strip, right=4pt of pol] (r1) {MFA fuerte};
  \node[terminal chip, strip, right=3pt of r1, terminal accent] (r2) {bloqueo 24 h tras cambio de SIM};
  \node[terminal chip, strip, right=3pt of r2] (r3) {todo canal cifrado};
  \node[terminal chip, strip, right=3pt of r3] (r4) {sesión 5 min, token 10 min};
\end{tikzpicture}

\end{frame}

% ------------------------------------------------------------------ 06
\section{Physical}
\begin{frame}{Physical \textmd{// mecanismos de seguridad}}
  \prompt{sabsa --layer physical --view builder --owner "Dir. Infraestructura"}
  \vskip0pt
  % diagrams/physical.tex — capa Physical: controles por anillos de defensa en profundidad,
% controles transversales y tiempos de recuperación (RTO/RPO).
\begin{tikzpicture}[x=1cm, y=1cm]
  % anillos de afuera hacia adentro: (gris de fondo, rótulo, controles)
  \foreach \i/\g/\capa/\ctl in {
      0/3/PROVEEDOR CLOUD/{certificaciones ISO 27001 y SOC 2, región alterna},
      1/7/PERÍMETRO/{WAF, anti-DDoS, bots, mTLS al sistema de pago},
      2/11/RED/{VPC por dominio, subredes privadas, microsegmentación},
      3/15/HOST Y CONTENEDORES/{imágenes firmadas, hardening CIS, runtime},
      4/19/APLICACIÓN/{gateway OAuth2, anti-replay, antifraude en línea, RASP}} {
    \pgfmathsetmacro{\T}{-0.47*\i}
    \pgfmathsetmacro{\B}{-4.6+0.31*\i}
    \pgfmathsetmacro{\L}{0.26*\i}
    \pgfmathsetmacro{\R}{9.3-0.26*\i}
    \draw[line width=0.7pt, fill=black!\g] (\L,\T) rectangle (\R,\B);
    \node[font=\tiny, anchor=north] at ({(\L+\R)/2},\T-0.05)
      {\textbf{\capa}\ \ \ctl};
  }
  % centro: los datos
  \node[fill=terminalaccent, text=white, font=\scriptsize\bfseries, align=center,
        minimum width=5.8cm, minimum height=1.0cm] (dat) at (4.65,-2.8)
    {DATOS\\[-1pt]{\mdseries\tiny firewall de BD, cifrado con llaves en HSM y tokenización}};

  % el atacante entra desde afuera y cada anillo lo frena
  \draw[terminal arrow, draw=terminalaccent, line width=1pt] (-0.6,-2.75) -- (0.95,-2.75);
  \node[font=\tiny\bfseries, text=terminalaccent, anchor=east, align=right] at (-0.05,-2.4) {ataque};
  \node[font=\tiny, anchor=north, text=black!70] at (4.65,-4.63)
    {si un anillo cae, el siguiente contiene el ataque};

  % controles transversales: atraviesan todos los anillos
  \begin{scope}[every node/.style={terminal box, anchor=north west, text width=3.3cm,
                                   align=left, font=\tiny, inner sep=2.5pt}]
    \node (idn) at (9.6,0.0)
      {\textbf{IDENTIDAD}\\passkeys para clientes, llave FIDO2 para empleados, PAM con acceso temporal};
    \node (mon) at ([yshift=-7pt]idn.south west)
      {\textbf{MONITOREO}\\SIEM 24/7, logs inmutables, alertas de fraude en vivo};
    \node (pro) at ([yshift=-7pt]mon.south west)
      {\textbf{PROCEDIMENTAL}\\segregación de funciones, doble aprobación, simulacros};
  \end{scope}
  \node[font=\tiny\bfseries, anchor=south west] at (9.6,0.02) {transversales};

  % tiempos de recuperación (RTO = cuánto puede estar caído; RPO = cuánto se puede perder)
  \node[terminal inv, anchor=north west, text width=3.3cm, align=left, font=\tiny\bfseries,
        inner sep=2.5pt] (th) at ([yshift=-9pt]pro.south west) {RTO / RPO};
  \node[draw, line width=0.6pt, anchor=north west, text width=3.3cm, align=left,
        font=\tiny, inner sep=2.5pt] at (th.south west)
    {\textcolor{terminalaccent}{\textbf{pagos: RPO 0}}\\
     RTO < 1 min (zona), ≤ 2 h (región)\\[1pt]
     llaves y QR: RTO 1 h, RPO 5 min\\[1pt]
     backoffice: RTO y RPO 24 h};
  % franja inferior: cada componente con su control (lo que el docente dijo que va a preguntar)
  \tikzset{map/.style={terminal chip, font=\tiny, text height=1.35ex, text depth=0.4ex, inner xsep=3pt}}
  \node[terminal tag, font=\tiny\bfseries, text height=1.35ex, text depth=0.4ex, anchor=west]
    (sm) at (0,-5.15) {COMPONENTE → CONTROL};
  \node[map, right=4pt of sm] (m1) {base de datos → firewall de BD y cifrado en HSM};
  \node[map, right=3pt of m1, terminal accent] (m2) {conexión Bre-B → canal dedicado, mTLS y firma ISO 20022};
  \node[map, anchor=north west] (m3) at ([yshift=-3pt]sm.south west) {clientes → passkeys (doble factor)};
  \node[map, right=3pt of m3] (m4) {pagos → antifraude en línea};
  \node[map, right=3pt of m4] (m5) {borde → WAF y anti-DDoS};
  \node[map, right=3pt of m5] (m6) {empleados → EDR y DLP};
\end{tikzpicture}

\end{frame}

% ------------------------------------------------------------------ 07
\section{Component}
\begin{frame}{Component \textmd{// productos y tecnologías}}
  \prompt{sabsa --layer component --view tradesman --owner "Líder SecOps"}
  \vskip4pt
  % diagrams/component.tex — capa Component: pila de productos (de internet hacia adentro, con la
% conexión a Bre-B en su propia banda) y agenda.
\begin{tikzpicture}[x=1cm, y=1cm]
  \tikzset{band/.style={draw=black, line width=0.5pt, minimum width=10.45cm, minimum height=0.56cm,
             anchor=north west, inner sep=0pt},
           lab/.style={terminal inv, text width=2.2cm, minimum height=0.56cm, anchor=north west,
             font=\scriptsize\bfseries, inner sep=0pt, align=center},
           prod/.style={terminal chip, text height=1.75ex, text depth=0.45ex}}
  % \capa{fila}{rótulo}{primer chip}{resto de chips separados por coma}
  \newcommand{\capa}[4]{%
    \node[band] (b#1) at (0,-0.66*#1) {};
    \node[lab] (l#1) at (b#1.north west) {#2};
    \node[prod, anchor=west] (c#1-0) at ([xshift=6pt]l#1.east) {#3};
    \foreach \t [count=\i from 1, remember=\i as \j (initially 0)] in {#4} {
      \node[prod, anchor=west] (c#1-\i) at ([xshift=4pt]c#1-\j.east) {\t};
    }}

  \node[font=\tiny, anchor=south west] at (0,0.05) {↓ internet, app y comercios};
  \capa{0}{PERÍMETRO}{Cloudflare (WAF, anti-DDoS, bots)}{API Gateway}
  \capa{1}{IDENTIDAD}{Auth0 / Keycloak}{Entra ID + YubiKey, CyberArk (PAM)}
  \capa{2}{PLATAFORMA}{EKS 3 AZ + DR}{Istio (mTLS), Aurora, MSK (Kafka)}
  \capa{3}{INTERBANCARIO}{Direct Connect (Equinix BG1)}{ISO 20022/REST, mTLS X.509}
  \capa{4}{CRIPTO}{CloudHSM}{KMS, HashiCorp Vault}
  \capa{5}{DETECCIÓN}{Feedzai}{Sentinel (SIEM), Purview DLP, CrowdStrike}
  \capa{6}{PIPELINE}{GitHub Actions}{Semgrep, Snyk, Trivy, ZAP}

  % en acento: la conexión a Bre-B (hilo de la integración segura) y el HSM (raíz de confianza)
  \node[lab, fill=terminalaccent] at (b3.north west) {INTERBANCARIO};
  \node[prod, terminal accent, line width=0.8pt, anchor=west] at (c4-0.west) {CloudHSM};

  % agenda de operación
  \node[terminal box, anchor=north west, text width=2.95cm, align=left, font=\scriptsize]
    (ag) at (10.75,-0.28) {%
      \textbf{>} pentest semestral\\\phantom{>} (lo exige la SFC)\\[2pt]
      \textbf{>} escaneo semanal\\\phantom{>} y en cada build\\[2pt]
      \textbf{>} DRP semestral\\[2pt]
      \textbf{>} tabletop anual\\\phantom{>} de ransomware\\[2pt]
      \textbf{>} Vault: rotación\\\phantom{>} diaria, certs 90 d\\[2pt]
      \textbf{>} chrony con la\\\phantom{>} Hora Legal del INM};
  \node[terminal inv, anchor=south west, text width=2.95cm, align=left] at (ag.north west) {agenda};
\end{tikzpicture}

\end{frame}

% ------------------------------------------------------------------ 08
\section{Conclusiones}
\begin{frame}{Conclusiones}
  \begin{cellitems}\small
    \item \textbf{El negocio manda:} El proyecto existe porque mueve un objetivo
      estratégico y la arquitectura justifica cada peso invertido.
    \vskip4pt
    \item \textbf{De arriba hacia abajo:} Bajamos desde lo que pide el
      dueño del negocio hasta los productos, de izquierda a derecha, y cada
      producto se puede rastrear hasta un requisito.
    \vskip4pt
    \item \textbf{La seguridad no se pierde al bajar de capa:} La integración con
      Bre-B pasa de confidencial, íntegra y disponible (negocio) a cifrado en tránsito, a mensajería segura,
      a mTLS con firma ISO 20022 y a Direct Connect, y entra al cronograma
      desde la compra de licencias.
    \vskip4pt
    \item \textbf{La seguridad es transversal:} No es un área de Tecnología sino
      un componente que toca a todas las áreas; por eso el CISO reporta a
      Presidencia.
    \vskip4pt
    \item \textbf{Lo que sigue:} La sexta capa de SABSA (operación) opera y mide
      esta arquitectura en el día a día.
  \end{cellitems}
  \vskip8pt
  \begin{ibox}Arquitectar = llevar el objetivo de negocio hasta los controles\end{ibox}
\end{frame}

\terminalscreen{\texttt{\$ sabsa -{}-preguntas}}{¿PREGUNTAS?}{\empresa{} \quad//\quad \proyecto}

\end{document}
